\begin{thebibliography}{99}

\bibitem{Bessenrodt}
Christine Bessenrodt.
Critical classes, Kronecker products of spin characters, and the Saxl conjecture.
\emph{Algebraic Combinatorics} \textbf{1} (2018), no.~3, 353--369.
\href{https://doi.org/10.5802/alco.18}{doi:10.5802/alco.18}.

\bibitem{BessenrodtBowman}
Christine Bessenrodt and Chris Bowman.
Splitting Kronecker squares, 2-decomposition numbers, Catalan combinatorics,
and the Saxl conjecture.
\emph{Algebraic Combinatorics} \textbf{6} (2023), no.~4, 863--899.
\href{https://doi.org/10.5802/alco.294}{doi:10.5802/alco.294}.

\bibitem{BessenrodtBowmanSutton}
Christine Bessenrodt, Chris Bowman, and Louise Sutton.
Kronecker positivity and 2-modular representation theory.
\emph{Transactions of the American Mathematical Society, Series B}
\textbf{8} (2021), 1024--1055.
\href{https://doi.org/10.1090/btran/70}{doi:10.1090/btran/70}.

\bibitem{BWZ}
Andrew A.~H. Brown, Stephanie van Willigenburg, and Mike Zabrocki.
Expressions for Catalan Kronecker products.
\emph{Pacific Journal of Mathematics} \textbf{248} (2010), no.~1, 31--48.
\href{https://doi.org/10.2140/pjm.2010.248.31}{doi:10.2140/pjm.2010.248.31}.

\bibitem{Ebrahimi}
Mahdi Ebrahimi.
The problem of deciding the positivity of Kronecker coefficients and Saxl conjecture.
Preprint, version~3, November 26, 2025.
\href{https://arxiv.org/abs/2511.03484v3}{arXiv:2511.03484v3}.

\bibitem{FultonHarris}
William Fulton and Joe Harris.
\emph{Representation Theory: A First Course}.
Graduate Texts in Mathematics, vol.~129.
Springer-Verlag, New York, 1991.
\href{https://doi.org/10.1007/978-1-4612-0979-9}{doi:10.1007/978-1-4612-0979-9}.

\bibitem{HarmanRyba}
Nate Harman and Christopher Ryba.
A tensor-cube version of the Saxl conjecture.
\emph{Algebraic Combinatorics} \textbf{6} (2023), no.~2, 507--511.
\href{https://doi.org/10.5802/alco.267}{doi:10.5802/alco.267}.

\bibitem{HSTZ}
Gerhard Heide, Jan Saxl, Pham Huu Tiep, and Alexandre E. Zalesski.
Conjugacy action, induced representations and the Steinberg square for
simple groups of Lie type.
\emph{Proceedings of the London Mathematical Society} (3)
\textbf{106} (2013), no.~4, 908--930.
\href{https://doi.org/10.1112/plms/pds062}{doi:10.1112/plms/pds062}.
Preprint: \href{https://arxiv.org/abs/1209.1768}{arXiv:1209.1768}.

\bibitem{Ikenmeyer}
Christian Ikenmeyer.
The Saxl conjecture and the dominance order.
\emph{Discrete Mathematics} \textbf{338} (2015), no.~11, 1970--1975.
\href{https://doi.org/10.1016/j.disc.2015.04.027}{doi:10.1016/j.disc.2015.04.027}.
Preprint: \href{https://arxiv.org/abs/1410.6549v2}{arXiv:1410.6549v2}.

\bibitem{James}
G.~D. James.
\emph{The Representation Theory of the Symmetric Groups}.
Lecture Notes in Mathematics, vol.~682.
Springer-Verlag, Berlin, 1978.
\href{https://doi.org/10.1007/BFb0067708}{doi:10.1007/BFb0067708}.

\bibitem{Letellier}
Emmanuel Letellier.
Tensor products of unipotent characters of general linear groups over finite fields.
\emph{Transformation Groups} \textbf{18} (2013), no.~1, 233--262.
\href{https://doi.org/10.1007/s00031-013-9211-3}{doi:10.1007/s00031-013-9211-3}.

\bibitem{LetellierNam}
Emmanuel Letellier and GyeongHyeon Nam.
The Saxl conjecture and the tensor square of unipotent characters of
$\GL_n(q)$.
\emph{Algebraic Combinatorics} \textbf{8} (2025), no.~4, 1119--1140.
\href{https://doi.org/10.5802/alco.434}{doi:10.5802/alco.434}.

\bibitem{Li}
Xin Li.
Saxl Conjecture for triple hooks.
\emph{Discrete Mathematics} \textbf{344} (2021), no.~6, article 112340.
\href{https://doi.org/10.1016/j.disc.2021.112340}{doi:10.1016/j.disc.2021.112340}.

\bibitem{LuoSellke}
Sammy Luo and Mark Sellke.
The Saxl conjecture for fourth powers via the semigroup property.
\emph{Journal of Algebraic Combinatorics} \textbf{45} (2017), no.~1, 33--80.
\href{https://arxiv.org/abs/1511.02387v2}{arXiv:1511.02387v2}.

\bibitem{Macdonald}
I.~G. Macdonald.
\emph{Symmetric Functions and Hall Polynomials}, second edition.
Oxford Mathematical Monographs.
Clarendon Press, Oxford, 1995.
\href{https://doi.org/10.1093/oso/9780198534891.001.0001}{doi:10.1093/oso/9780198534891.001.0001}.

\bibitem{SaxlCyclic}
{\raggedright
OpenAI.
A Cyclic Polytabloid Proof of Saxl's Conjecture.
OpenAI Math Release preprint
\href{https://github.com/openai/math/blob/main/preprints/A-Cyclic-Polytabloid-Proof-of-Saxls-Conjecture-September-24-2026/paper.pdf}{OAI:A-Cyclic-Polytabloid-Proof-of-Saxls-Conjecture-September-24-2026},
2026. Theorem~3.1, proved in Section~6.
\par}

\bibitem{PPV}
Igor Pak, Greta Panova, and Ernesto Vallejo.
Kronecker products, characters, partitions, and the tensor square conjectures.
\emph{Advances in Mathematics} \textbf{288} (2016), 702--731.
\href{https://doi.org/10.1016/j.aim.2015.11.002}{doi:10.1016/j.aim.2015.11.002}.
Preprint: \href{https://arxiv.org/abs/1304.0738v2}{arXiv:1304.0738v2}.
Cited numbering follows this version.

\bibitem{Zelevinsky}
Andrei Zelevinsky.
Littlewood--Richardson semigroups.
In \emph{New Perspectives in Algebraic Combinatorics},
Louis J. Billera, Anders Bj{\"o}rner, Curtis Greene, Rodica Simion,
and Richard P. Stanley (eds.),
Mathematical Sciences Research Institute Publications, vol.~38,
Cambridge University Press, Cambridge, 1999, pp.~337--345.
\href{https://library.slmath.org/books/Book38/files/zelevinsky.pdf}{Institutional full text}.

\end{thebibliography}
